%% kapa: exam=1 round=32 year=2021 period=2 form=A subject=appraisal_law subject_ko=감정평가관계법규 seq=14 number=14
%% source: Q-Net 감정평가사 기출문제 (cst00309) artlSeq=5208397 file=extracted/1008/5208397/A형/2021년도 제32회 감정평가사 1차 2교시 A형.pdf
%% conversion: kordoc 4.24.51 -> tools/gen1.py (auto); review=proofread-sonnet; flags=none
\begin{question}{14}
부동산 가격공시에 관한 법령상 개별공시지가에 관한 설명으로 옳지 않은 것은?
\begin{choices}
  \item 시장·군수 또는 구청장은 개별공시지가에 토지가격비준표의 적용에 오류가 있음을 발견한 때에는 지체 없이 이를 정정하여야 한다.
  \item 표준지로 선정된 토지에 대하여 개별공시지가를 결정·공시하지 아니하는 경우에는 해당 토지의 표준지공시지가를 개별공시지가로 본다.
  \item 개별공시지가에 이의가 있는 자는 그 결정·공시일부터 60일 이내에 서면 또는 구두로 이의를 신청할 수 있다.
  \item 개별공시지가의 결정·공시에 소요되는 비용 중 국고에서 보조할 수 있는 비용은 개별공시지가의 결정·공시에 드는 비용의 50퍼센트 이내로 한다.
  \item 개별공시지가의 단위면적은 1제곱미터로 한다.
\end{choices}
\end{question}
\answer{3}
